\section{The PostPC and the new computational requirements}

Процесът на трансформация на формите на човешкия труд поставя качествено нови проблеми за решаване от едно общество. 
Компютърните технологии като незаменим инструмент в решаването на по - голямата част от тези проблеми, също трябва 
да се адаптират подходящо, за да откликнат на нуждите / потребностите на новите стопански и социални реалности.
Нещо повече, оказва се, че промените в компютърните технологии, подобно на социално - икономическите изменения следват
принципите на еволюцията и ко-еволюцията. И наистина, първите разработени компютри са програмирани на машинен код,
по - късно, за улеснение машинните инструкции били заменени с мнемонични, които са тяхно биективно съответствие.
Машинния код и мнемоничните езици като семантика на инструкциите са много по - близко до програмирането на 
физическата памет и хардуера на един компютър и в този смисъл те не са удобни за програмирането на абстрактни процеси. 
Това неудобство е породило нуждата да се изгради една по - абстрактна архитектура от инструкции (\acr{ISA}), така, че
да се описват процеси, които да са от по - широк научен интерес. Тъй като първите научни кръгове, в които компютърните 
технологии са станали достояние са били тези на математиците и инженерите, то естествено е да се очаква, че първите
езици за програмиране след мнемоничните ще имат \acr{ISA}, което е изключително удобно за разработката и програмирането на
алгоритми и математически изчисления. Така първите езици като \acr{Fortran} (Formula translator) и \acr{Algol} били създадени предимно
за изграждането на алгоритми и математически изчисления. Те използвали близък до математическия език \acr{ISA}, който 
посредством разбор се "превеждал" или както е утвърдено да се нарича "компилирал" до асемблер, а от там до машинен код. 
Оказало се обаче, че тези езици са не дотам удобни при програмирането на задачи от стопански и производствен характер 
(писане на операционна система, правене на програма за складови наличности, четене на данни от много входове и др).
Търсенията на програмистите за удачен език след този период съвсем не следвали някакви закономерности, а по - скоро 
имали случайностен характер. Например езика Pascal бил разработен с цел да служи за преподаване на студентите, 
\acr{Ada} бил разработен <опиши защо, понеже не мога да се сетя>, basic бил разработен (опиши защо). Най - сполучлив се 
оказал езика за програмиране C, създаден от (напиши името) през (напиши годината).Този език притежавал най - добрата 
комбинация от \acr{ISA}, които са удобни и за алгоритмичен и за хардуерен дизаин. С него се пишело удобно както алгоритмично 
поставен проблем, така и програмен проблем (като например писането на език. Даже компилатора на C е написан на самия език C). 
Освен това, той бил и достатъчно ефективен откъм бързодействие. Тук пак компилатора изиграл много важна роля, 
защото той имал за задача да сведе инструкциите на C до инструкции на асемблер. Благдарение на езика C се създали
много и различни по вид езици като например Prolog <назови още>. Някои от езиците приличали по синтаксис 
и функционалности на C / C++, а други се стремели да се разграничат от неговата семантика. Заслужава внимание езика Java,
който е подобен на C++, но кодът се компилира не до машинен, а до JVC, който след това се превежда до машинен 
за всяка една компютърна архитектура, така че да е максимално ефективен. Освен това при Java не се допуска менажиране на паметта. 
Други езици се отдалечили от парадигмата на C, като JavaScript и Python например. Те напълно премахнали
типизирането на данните и се съсредоточили върху \acr{ISA}, който да е оптимален за <опиши с няколко думи>. 
Постепено в практиката се наложили езиците, които не използват типове от данни и са с опростен синтаксис и не съдържат
или съдържат ограничен брой машинни инструкции. Практиката и характера на програмните продукти утвърдила парадигмата 
на обектно ориентираното програмиране пред останалите. Ако пътят на програмирането беше праволинеен 
или подвластен на някакъв закон би се очаквало тези тенденции да се затвърждават с течението на времето. 
Случва се обаче тозно обратното. След утвърждаването на уеб технологиите, потребността от по - ефективо изпълнение на
даден процес отново връща на мода строго типизираните езици, като най - разпространени стават Java и C#.
JavaScript откликва на това предизвикателство като се превръжа и в backend език чрез технологията NodeJS 
(по - късно и чрез Deno и Bun) и чрез "псевдотипизирането" си в езика TypeScript. Python откликва чрез използването 
на \acr{LLVM} интерпретатори като Numba или чрез нови компилатори cpy, като също вевежда псевдо типизиране. 
Но историята съвсем не свършва до тук. Писането на все по - ефективни програми, обработващи големи данни,
става все по - голямо предизвикателство и едно от най - простото работещо решение е паралелната обработка. 
Отначало паралелната обработка се свеждала до нишково програмиране, в последствие се разрастнало до 
паралелно свързани изчислителни машини, а в днешно време се е утвърдила парадигмата на хетерогенното програмиране,
с която изпълнението на програмите се разпределя между централния процесор и графичния процесор. 
Това доведе отново до „възцаряването“ на езика C, който поддържа програмирането на графична карта и протокол за 
управление на паметта на графичната карта, а строготипизираните данни са задължително условие 
при изпълняването на код върху графичната карта. Езиците като Python намират решение на този проблем като 
създават компилатор хетерогенни програми с помощта на \acr{LLVM}, а други като TypeScript, чрез \acr{FFI} протокол.
За отбелязване е, че всички езици ползват \acr{FFI} протокол, понеже само C може да управлява паметта на графичната карта.


Внимателният читател без труд ще забележи, че и в този случай
развитието на езиците за програмиране не следва някакъв строго определен праволинеен ход, а е низ от случайности,
потребности на обществото и резерв на изменчивост, който езиците притежават, за да откликнат на нуждите на 
нововъзникващата среда. Така хода на развитието на програмните езици има по - скоро зигзакообразен вид и 
е еволюционен по своята природа, а не детерминистичен. Друго достойно за отбелязване явление е, че отново възниква
нуждата от компилатори, чрез които да се оптимизира ефективостта на кода. Последните две твърдения са били открити и
изложени от крупния научен работник, аристократ и общественик Петър Свещаров. Благодарение на него аз се вдъхнових 
и изградих този проект.

